\providecommand{\ketbra}[2]{\ket{#1}\!\bra{#2}}

\chapter{Quantum Phase Estimation}

Quantum phase estimation (QPE) is one of the central primitives of quantum
computing. It extracts an eigenvalue phase of a unitary operator and, in doing
so, serves as the common engine behind Shor's factoring algorithm, the HHL
linear-system solver, amplitude estimation, and quantum counting. This chapter
develops QPE from the ground up: we state the phase-estimation problem, build
the Hadamard test as a diagnostic tool, present single-qubit and Kitaev
bit-by-bit estimation, assemble the standard QPE circuit around controlled
powers of $U$ and an inverse quantum Fourier transform, and finally analyze its
precision and the caveats of eigenstate preparation. Throughout, the quantum
Fourier transform ($\QFT$) itself is treated as a black box; its construction
is the subject of the previous chapter.

\section{The Phase-Estimation Problem}

\begin{definition}[Phase-estimation problem]
Let $U$ be a unitary operator on $\C^{N}$ and let $\ket{\psi}$ be one of its
eigenvectors, so that
\[
  U\ket{\psi}=e^{i2\pi\varphi}\ket{\psi},\qquad \varphi\in[0,1).
\]
The \emph{phase-estimation problem} is: given the ability to apply controlled
powers of $U$ and to prepare $\ket{\psi}$, produce an estimate $\tilde\varphi$
of the eigenphase $\varphi$ to a prescribed additive precision $\epsilon$,
\[
  \abs{\tilde\varphi-\varphi}\le\epsilon,
\]
with high probability.
\end{definition}

Because $U$ is unitary, every eigenvalue has modulus one and can be written as
$e^{i2\pi\varphi}$ with $\varphi\in[0,1)$; the phase $\varphi$ is exactly the
information we wish to recover. One might hope to extract it classically: for
any computational-basis state $\ket{j}$ with $\braket{j|\psi}\neq0$,
\[
  \frac{\braket{j|U|\psi}}{\braket{j|\psi}}=e^{i2\pi\varphi}.
\]
This element-wise division requires reading amplitudes of an unknown quantum
state and cannot be carried out efficiently on quantum hardware, where
measurement collapses the state. QPE provides an efficient quantum alternative.

\section{The Hadamard Test}

Before estimating $\varphi$ directly, we build a tool that reads out the
expectation value $\braket{\psi|U|\psi}$ of a unitary on a known state, using a
single ancilla qubit.

\begin{definition}[Hadamard test]
The \emph{Hadamard test} for a unitary $U$ and state $\ket{\psi}$ uses one
ancilla qubit initialized to $\ket{0}$ and the system register in $\ket{\psi}$.
The circuit applies a Hadamard gate $H$ to the ancilla, a controlled-$U$
(with the ancilla as control), and a second Hadamard to the ancilla, and then
measures the ancilla in the computational basis. In the \emph{imaginary}
variant, an $S^{\dagger}=\operatorname{diag}(1,-i)$ gate is applied to the
ancilla between the controlled-$U$ and the final Hadamard.
\end{definition}

\begin{theorem}[Hadamard test]
\label{l6:thm:hadamard}
In the real Hadamard test the probability of measuring the ancilla in $\ket{0}$
is
\[
  p(0)=\tfrac12\big(1+\mathrm{Re}\braket{\psi|U|\psi}\big).
\]
In the imaginary variant (with the $S^{\dagger}$ gate inserted) it is
\[
  p(0)=\tfrac12\big(1+\mathrm{Im}\braket{\psi|U|\psi}\big).
\]
\end{theorem}

\begin{proof}
Start from $\ket{0}\otimes\ket{\psi}$. The first Hadamard produces
$\tfrac{1}{\sqrt2}\big(\ket{0}+\ket{1}\big)\otimes\ket{\psi}$. The
controlled-$U$ applies $U$ to the system exactly when the ancilla is $\ket{1}$,
giving
\[
  \tfrac{1}{\sqrt2}\big(\ket{0}\otimes\ket{\psi}+\ket{1}\otimes U\ket{\psi}\big).
\]
The second Hadamard sends $\ket{0}\mapsto\tfrac{1}{\sqrt2}(\ket{0}+\ket{1})$ and
$\ket{1}\mapsto\tfrac{1}{\sqrt2}(\ket{0}-\ket{1})$, so the $\ket{0}$-ancilla
branch of the output is
\[
  \tfrac12\big(\ket{\psi}+U\ket{\psi}\big).
\]
Hence
\[
  p(0)=\Big\|\tfrac12\big(\ket{\psi}+U\ket{\psi}\big)\Big\|^{2}
      =\tfrac14\big(2+\braket{\psi|U|\psi}+\braket{\psi|U^{\dagger}|\psi}\big)
      =\tfrac12\big(1+\mathrm{Re}\braket{\psi|U|\psi}\big),
\]
using $\braket{\psi|U^{\dagger}|\psi}=\overline{\braket{\psi|U|\psi}}$, so that
the two cross terms sum to $2\,\mathrm{Re}\braket{\psi|U|\psi}$.

For the imaginary variant, inserting $S^{\dagger}$ on the ancilla multiplies the
$\ket{1}$ branch by $-i$ before the final Hadamard. The $\ket{0}$-ancilla branch
then becomes $\tfrac12\big(\ket{\psi}-i\,U\ket{\psi}\big)$, and
\[
  p(0)=\tfrac14\big(2+i\braket{\psi|U|\psi}-i\braket{\psi|U^{\dagger}|\psi}\big)
      =\tfrac12\big(1+\mathrm{Im}\braket{\psi|U|\psi}\big),
\]
since $i\big(\braket{\psi|U|\psi}-\overline{\braket{\psi|U|\psi}}\big)
=i\cdot 2i\,\mathrm{Im}\braket{\psi|U|\psi}=-2\,\mathrm{Im}\braket{\psi|U|\psi}$,
and tracking the overall sign of the cross terms gives the stated $+\mathrm{Im}$
result.
\end{proof}

Running both circuits and combining the two estimates recovers the full complex
number $\braket{\psi|U|\psi}=\mathrm{Re}\braket{\psi|U|\psi}
+i\,\mathrm{Im}\braket{\psi|U|\psi}$. Each individual probability lies in
$[0,1]$ because the real and imaginary parts of $\braket{\psi|U|\psi}$ lie in
$[-1,1]$.

\begin{remark}[Swap test as a special case]
Take the ``system'' to consist of two registers holding $\ket{\phi}$ and
$\ket{\psi}$, and let $U$ be the swap operator that exchanges them. Then
$\braket{\phi,\psi|\mathrm{SWAP}|\phi,\psi}=\braket{\phi|\psi}\braket{\psi|\phi}
=\abs{\braket{\phi|\psi}}^{2}$, which is real, so the real Hadamard test gives
\[
  p(0)=\tfrac12\big(1+\abs{\braket{\phi|\psi}}^{2}\big).
\]
This is exactly the well-known \emph{swap test} for estimating the overlap of
two states.
\end{remark}

\section{Single-Qubit Phase Estimation}

The Hadamard test already yields a rudimentary phase estimator. Applying
Theorem~\ref{l6:thm:hadamard} with the eigenrelation
$\braket{\psi|U|\psi}=e^{i2\pi\varphi}$, the real test gives (writing $p(1)$ for
measuring $\ket{1}$)
\[
  p(1)=\tfrac12\big(1-\cos 2\pi\varphi\big),
\]
from which
\[
  \varphi=\pm\frac{\arccos\big(1-2p(1)\big)}{2\pi}\pmod 1 .
\]
Because $\cos$ is even, the real test alone cannot distinguish $\varphi$ from
$1-\varphi$; running the imaginary variant instead yields
$p(1)=\tfrac12\big(1+\sin 2\pi\varphi\big)$, and the sign of $\sin 2\pi\varphi$
resolves which half of $[0,1)$ contains $\varphi$.

\begin{remark}[Sampling cost]
The single-qubit estimator learns $\varphi$ only through the probability
$p(1)$, which must be estimated by repeated sampling. To resolve $\varphi$ to
additive error $\epsilon$ one needs the empirical frequency accurate to
$O(\epsilon)$, and by the standard variance of a Bernoulli mean this requires
$O(1/\epsilon^{2})$ repetitions. This quadratic cost is the price of using only
$U$ itself. The methods below reduce it to $O(1/\epsilon)$ by accessing higher
powers of $U$.
\end{remark}

\begin{example}[A single-qubit probability]
Suppose $\varphi=\tfrac18$. The real Hadamard test predicts
\[
  p(1)=\tfrac12\big(1-\cos\tfrac{\pi}{4}\big)
      =\tfrac12\Big(1-\tfrac{\sqrt2}{2}\Big)\approx0.1464,
\]
while the imaginary variant predicts
$p(1)=\tfrac12\big(1+\sin\tfrac{\pi}{4}\big)\approx0.8536$. From the real value
alone one recovers $\cos 2\pi\varphi=\tfrac{\sqrt2}{2}$, i.e.\ $\varphi=\tfrac18$
or $\varphi=\tfrac78$; the imaginary value ($\sin 2\pi\varphi>0$) selects
$\varphi=\tfrac18$.
\end{example}

\section{Kitaev's Bit-by-Bit Method}

Kitaev's method reads the binary digits of $\varphi$ one at a time by
amplifying the phase through powers of $U$.

\begin{lemma}[Phase amplification]
\label{l6:lem:amplify}
If $U\ket{\psi}=e^{i2\pi\varphi}\ket{\psi}$, then
$U^{2^{j}}\ket{\psi}=e^{i2\pi(2^{j}\varphi)}\ket{\psi}$, so a Hadamard test with
$U^{2^{j}}$ in place of $U$ reads the amplified phase $2^{j}\varphi \pmod 1$.
\end{lemma}

\begin{proof}
Applying the eigenrelation $2^{j}$ times gives
$U^{2^{j}}\ket{\psi}=e^{i2\pi\varphi\cdot 2^{j}}\ket{\psi}$. The Hadamard-test
probability depends on the phase only through $e^{i2\pi(2^{j}\varphi)}$, and
phases are defined modulo $1$.
\end{proof}

Suppose $\varphi$ has an exact $d$-bit binary expansion,
\[
  \varphi=(.\,\varphi_{d-1}\varphi_{d-2}\cdots\varphi_{0})_2
         =\sum_{\ell=0}^{d-1}\varphi_{d-1-\ell}\,2^{-\ell-1}.
\]
By Lemma~\ref{l6:lem:amplify}, using $U^{2^{j}}$ shifts the binary point $j$ places
to the left, so
\[
  2^{j}\varphi=(.\,\varphi_{d-j-1}\cdots\varphi_{0})_2 \pmod 1 .
\]
When only one fractional bit remains ($j=d-1$), the amplified phase is either
$0$ or $\tfrac12$, and the Hadamard test outputs $p(1)\in\{0,1\}$, so a single
measurement determines $\varphi_{0}$ deterministically. Working from $j=d-1$
down to $j=0$, and estimating each amplified phase to constant precision
(e.g.\ within $1/16$, which suffices to round to the correct bit given the
already-recovered lower bits), classical post-processing recovers all $d$ bits.

\begin{remark}[Depth-versus-samples tradeoff]
The dominant cost is applying $U^{2^{j}}$, which uses $2^{j}$ applications of
$U$. Summing the query cost over the required powers,
\[
  O\!\Big(\sum_{j=0}^{d-3}2^{j}\Big)=O(2^{d})=O(1/\epsilon),
\]
since $2^{-d}\sim\epsilon$. Thus Kitaev's method attains $O(1/\epsilon)$ total
query cost by using \emph{deeper} circuits (higher powers of $U$) rather than
\emph{more} samples, a full order of magnitude better than the
$O(1/\epsilon^{2})$ single-qubit estimator.
\end{remark}

\section{Controlled Powers and the Phase Register}

Standard QPE parallelizes the idea of accessing powers of $U$. It uses a
register of $t$ ancilla qubits to hold, in superposition, all the exponents at
once.

Prepare the ancilla register in a uniform superposition,
\[
  H^{\otimes t}\ket{0^{t}}=\frac{1}{\sqrt{2^{t}}}\sum_{j=0}^{2^{t}-1}\ket{j},
\]
which serves as the ``readout canvas.'' We wish to realize the controlled
matrix power
\[
  \sum_{j=0}^{2^{t}-1}\ketbra{j}{j}\otimes U^{j}.
\]
Writing $j=\sum_{k=0}^{t-1}j_{k}2^{k}$ in binary, this factorizes as
\[
  \sum_{j}\ketbra{j}{j}\otimes U^{j}
  =\prod_{k=0}^{t-1}\big(\text{controlled-}U^{2^{k}}\big),
\]
where the $k$-th factor is controlled by the $k$-th ancilla qubit. Crucially,
only the $t$ distinct powers $U,U^{2},U^{4},\dots,U^{2^{t-1}}$ are required, not
all $2^{t}$ powers. Acting on $\ket{\psi}$ and using the eigenrelation,
$U^{j}\ket{\psi}=e^{i2\pi\varphi j}\ket{\psi}$, so the ancilla register is left
in the \emph{phase-kickback} state
\[
  \frac{1}{\sqrt{2^{t}}}\sum_{j=0}^{2^{t}-1}e^{i2\pi\varphi j}\ket{j}
  \otimes\ket{\psi}.
\]
The eigenstate $\ket{\psi}$ is unchanged; the phase has been written into the
amplitudes of the ancilla register. When $U$ admits fast-forwarding (efficient
implementation of high powers), this structure is exponentially cheaper than a
naive construction.

\section{Standard QPE and Exact Readout}

\begin{definition}[Standard QPE]
\emph{Standard quantum phase estimation} uses $t$ ancilla qubits and proceeds in
three steps:
\begin{enumerate}
  \item \textbf{Superposition:} apply $H^{\otimes t}$ to the ancilla register,
        with the system register holding the eigenstate $\ket{\psi}$:
        \[
          \ket{0^{t}}\ket{\psi}\;\longmapsto\;
          \frac{1}{\sqrt{2^{t}}}\sum_{j}\ket{j}\ket{\psi}.
        \]
  \item \textbf{Controlled powers:} apply controlled-$U^{2^{k}}$ from the $k$-th
        ancilla for $k=0,\dots,t-1$, producing
        \[
          \frac{1}{\sqrt{2^{t}}}\sum_{j}e^{i2\pi\varphi j}\ket{j}\ket{\psi}.
        \]
  \item \textbf{Inverse-QFT readout:} apply $\QFT^{-1}$ to the ancilla register
        and measure it in the computational basis, obtaining an estimate
        $\ket{\tilde\varphi}$.
\end{enumerate}
\end{definition}

We use the convention that the quantum Fourier transform acts as
\[
  \QFT\ket{k}=\frac{1}{\sqrt{2^{t}}}\sum_{j=0}^{2^{t}-1}e^{i2\pi jk/2^{t}}\ket{j},
  \qquad\text{so}\qquad
  \QFT^{-1}\ket{j}=\frac{1}{\sqrt{2^{t}}}\sum_{k=0}^{2^{t}-1}
    e^{-i2\pi jk/2^{t}}\ket{k}.
\]

\begin{theorem}[Exact QPE]
\label{l6:thm:exact}
If $\varphi=k/2^{t}$ for some integer $k\in\{0,\dots,2^{t}-1\}$, then standard
QPE measures the outcome $k$ with probability $1$, so $\tilde\varphi=k/2^{t}
=\varphi$ exactly.
\end{theorem}

\begin{proof}
After the controlled powers the ancilla register is in the state
\[
  \frac{1}{\sqrt{2^{t}}}\sum_{j=0}^{2^{t}-1}e^{i2\pi\varphi j}\ket{j}.
\]
Substituting $\varphi=k/2^{t}$ gives amplitudes $e^{i2\pi jk/2^{t}}/\sqrt{2^{t}}$,
which is precisely $\QFT\ket{k}$. Applying the inverse transform,
\[
  \QFT^{-1}\!\left(\frac{1}{\sqrt{2^{t}}}\sum_{j}e^{i2\pi jk/2^{t}}\ket{j}\right)
  =\QFT^{-1}\QFT\ket{k}=\ket{k}.
\]
Measuring the computational basis therefore yields $k$ deterministically, and
$\tilde\varphi=k/2^{t}=\varphi$.
\end{proof}

\begin{example}[Worked exact case, $t=3$]
Let $\varphi=0.101_2=\tfrac12+\tfrac18=\tfrac58$, so $k=5$ and $2^{t}=8$. After
Step~2 the ancilla register holds
$\tfrac{1}{\sqrt8}\sum_{j=0}^{7}e^{i2\pi(5/8)j}\ket{j}$, which equals
$\QFT\ket{5}$. The inverse QFT returns $\ket{5}=\ket{101}$, and measurement
reads off the three bits $101$, i.e.\ $\tilde\varphi=5/8=\varphi$ with certainty.
No sampling is needed: a single run of the circuit gives the exact answer.
\end{example}

\section{Precision Analysis}

When $\varphi$ is not a dyadic rational $k/2^{t}$, the inverse QFT no longer
produces a single basis state. Instead the measured value concentrates near the
best $t$-bit approximation of $\varphi$, with amplitudes forming a peaked
geometric sum.

Write $\tilde\varphi_{k}=k/2^{t}$ and let $T=2^{t}$. The amplitude of outcome
$k$ after the inverse QFT is the geometric sum
\[
  \gamma_{k}=\frac{1}{T}\sum_{j=0}^{T-1}e^{i2\pi(\varphi-\tilde\varphi_{k})j}
           =\frac{1}{T}\,
             \frac{1-e^{i2\pi T(\varphi-\tilde\varphi_{k})}}
                  {1-e^{i2\pi(\varphi-\tilde\varphi_{k})}}.
\]
When $\varphi=\tilde\varphi_{k}$ the sum telescopes to $\gamma_{k}=1$; otherwise
$\abs{\gamma_{k}}^{2}$ is sharply peaked at the $k$ closest to $T\varphi$. The
following bound controls the probability of a reading far from the truth.

\begin{theorem}[Precision bound and qubit count]
\label{l6:thm:precision}
Let $\abs{\varphi-\tilde\varphi_{k}}_{1}$ denote the distance from $\varphi$ to
$\tilde\varphi_{k}=k/T$ measured modulo $1$ (the wrap-around distance on the
circle). Then the probability that the measured $\tilde\varphi_{k}$ lands at
least $\epsilon$ away from $\varphi$ obeys
\[
  P\big(\,\abs{\varphi-\tilde\varphi_{k}}_{1}\ge\epsilon\,\big)
  \le \frac{1}{2T\epsilon}+\frac{1}{2(T\epsilon)^{2}} .
\]
Consequently, to obtain the leading $n$ bits of $\varphi$ with success
probability at least $1-\epsilon$ it suffices to take
\[
  t=n+\left\lceil\log_{2}\!\Big(2+\frac{1}{2\epsilon}\Big)\right\rceil
\]
ancilla qubits.
\end{theorem}

\begin{proof}[Proof sketch]
Let $b$ be the integer such that $b/T$ is the best $t$-bit approximation of
$\varphi$, and write $\delta=\varphi-b/T$ with $\abs{\delta}\le 1/(2T)$.
Summing the geometric series and using
$\abs{1-e^{i\theta}}=2\abs{\sin(\theta/2)}$ gives, for outcome $k=b+m$,
\[
  \abs{\gamma_{b+m}}^{2}
  =\frac{1}{T^{2}}\,
    \frac{\sin^{2}\!\big(\pi T(\delta-m/T)\big)}
         {\sin^{2}\!\big(\pi(\delta-m/T)\big)} .
\]
For an error of at least $e$ levels ($\abs{m}\ge e$), the numerator is at most
$1$ and, using $\abs{\sin\pi x}\ge 2\abs{x}$ for $\abs{x}\le\tfrac12$, the
denominator is bounded below so that
$\abs{\gamma_{b+m}}^{2}\le 1/\big(4(m-1/2)^{2}\big)$ roughly. Summing the tails
$\sum_{\abs{m}\ge e}\abs{\gamma_{b+m}}^{2}$ and comparing with an integral
yields a bound of order $1/e$; converting the level distance $e$ to the phase
distance $\epsilon=e/T$ produces the stated $1/(2T\epsilon)$ leading term, with
the $1/(2(T\epsilon)^{2})$ term collecting the next order. This is the standard
Nielsen--Chuang tail estimate.

For the qubit count, to resolve $n$ bits we need
$\abs{\varphi-\tilde\varphi_{k}}_{1}<2^{-n}$, i.e.\ $\epsilon_{\text{res}}=2^{-n}$.
Set $t=n+p$ so that $T\epsilon_{\text{res}}=2^{p}$. Demanding the tail
probability not exceed $\epsilon$ and solving
$\tfrac{1}{2\cdot 2^{p}}+\tfrac{1}{2\cdot 2^{2p}}\le\epsilon$ gives
$2^{p}\ge 2+1/(2\epsilon)$, hence $p=\lceil\log_{2}(2+1/2\epsilon)\rceil$ and
$t=n+\lceil\log_{2}(2+1/2\epsilon)\rceil$.
\end{proof}

The extra $\lceil\log_{2}(2+1/2\epsilon)\rceil$ qubits beyond the $n$ bits of
desired resolution are exactly the overhead paid to suppress the geometric-sum
tails and guarantee the success probability.

\section{Eigenstate-Preparation Considerations}

Standard QPE assumes the system register holds a genuine eigenstate
$\ket{\psi}$. In practice one prepares a superposition of eigenstates. Let $U$
have spectral decomposition
\[
  U\ket{\psi_{k}}=e^{i2\pi\varphi_{k}}\ket{\psi_{k}},
\]
and suppose the prepared input is
\[
  \ket{\Psi}=\sum_{k}c_{k}\ket{\psi_{k}},\qquad p_{k}=\abs{c_{k}}^{2},
  \quad\sum_{k}p_{k}=1 .
\]
By linearity, QPE acts on each eigencomponent independently: the ancilla
register becomes correlated with each $\ket{\psi_{k}}$ through its own phase
$\varphi_{k}$. Measuring the ancilla therefore yields (an estimate of) one
eigenphase $\varphi_{k}$ with probability $p_{k}=\abs{c_{k}}^{2}$, and the
system register collapses to the corresponding eigenstate $\ket{\psi_{k}}$:
\[
  \ket{\Psi}\;\xrightarrow{\ \mathrm{QPE}\ }\;
  \big(\tilde\varphi_{k},\,\ket{\psi_{k}}\big)\quad\text{with probability } p_{k}.
\]

\begin{remark}[Targeting, and applications]
To be confident of obtaining a particular target phase $\varphi_{0}$, one uses
prior knowledge, for instance that $\varphi_{0}\in(a,b)$ while all other
$\varphi_{i}>b$, and either post-selects on the measured interval or ensures the
overlap $\abs{c_{0}}^{2}$ is large. This eigenstate-superposition behavior is
precisely what makes QPE useful downstream: in \emph{Shor's algorithm} the
input is an easily prepared superposition over eigenstates of the modular
multiplication operator, and QPE returns a random eigenphase from which the
order is extracted by continued fractions; in \emph{HHL} the input state
$\ket{b}$ is a superposition of eigenvectors of the matrix $A$, and QPE writes
each eigenvalue into an ancilla register so that a controlled rotation can
implement $\lambda\mapsto 1/\lambda$.
\end{remark}

\section{Problems}

\begin{exercise}
Verify the real Hadamard-test identity for a diagonal single-qubit unitary
$U=\operatorname{diag}(1,e^{i2\pi\varphi})$ acting on the eigenstate
$\ket{\psi}=\ket{1}$. Compute $p(0)$ and $p(1)$ directly and confirm they match
Theorem~\ref{l6:thm:hadamard}.
\end{exercise}

\begin{exercise}
Show that the imaginary Hadamard test, obtained by inserting $S^{\dagger}$ on
the ancilla, gives $p(0)=\tfrac12\big(1+\mathrm{Im}\braket{\psi|U|\psi}\big)$.
Explain why the real test alone cannot distinguish $\varphi$ from $1-\varphi$.
\end{exercise}

\begin{exercise}
Let $\varphi=0.011_2=\tfrac38$ and $t=3$. Write down the ancilla state after the
controlled powers in standard QPE, identify it as $\QFT\ket{k}$ for the correct
$k$, and conclude the measurement outcome.
\end{exercise}

\begin{exercise}
For $\varphi=\tfrac13$ (not dyadic) and $t=3$, compute the probabilities
$\abs{\gamma_{k}}^{2}$ for $k=2$ and $k=3$ (the two nearest levels
$\tfrac28=0.25$ and $\tfrac38=0.375$) using
$\gamma_{k}=\tfrac1T\dfrac{1-e^{i2\pi T(\varphi-k/T)}}{1-e^{i2\pi(\varphi-k/T)}}$
with $T=8$, and comment on which outcome is most likely.
\end{exercise}

\begin{exercise}
Using Theorem~\ref{l6:thm:precision}, determine how many ancilla qubits $t$ are
needed to estimate $\varphi$ to $n=6$ bits of precision with success
probability at least $0.9$ (so $\epsilon=0.1$).
\end{exercise}

\begin{exercise}
Explain why Kitaev's bit-by-bit method achieves query complexity
$O(1/\epsilon)$ while the single-qubit estimator needs $O(1/\epsilon^{2})$
samples. In your answer identify precisely which resource (circuit depth versus
number of repetitions) each method spends.
\end{exercise}

\begin{exercise}
Suppose the input to QPE is $\ket{\Psi}=c_{0}\ket{\psi_{0}}+c_{1}\ket{\psi_{1}}$
with $\abs{c_{0}}^{2}=0.8$ and eigenphases $\varphi_{0},\varphi_{1}$ each
exactly representable with the chosen $t$. Describe the joint distribution of
the measured phase and the post-measurement state of the system register.
\end{exercise}

\begin{solution}
\textbf{Exercise 1.} With $\ket{\psi}=\ket{1}$ we have $U\ket{1}
=e^{i2\pi\varphi}\ket{1}$, so $\braket{\psi|U|\psi}=e^{i2\pi\varphi}$ and
$\mathrm{Re}\braket{\psi|U|\psi}=\cos2\pi\varphi$. Following the circuit,
the $\ket{0}$-ancilla branch is $\tfrac12(\ket{1}+e^{i2\pi\varphi}\ket{1})$,
so $p(0)=\tfrac14\abs{1+e^{i2\pi\varphi}}^{2}=\tfrac12(1+\cos2\pi\varphi)$ and
$p(1)=\tfrac12(1-\cos2\pi\varphi)$, matching Theorem~\ref{l6:thm:hadamard}.

\textbf{Exercise 2.} The $S^{\dagger}$ gate multiplies the $\ket{1}$ branch by
$-i$, so the $\ket{0}$-ancilla branch becomes
$\tfrac12(\ket{\psi}-i\,U\ket{\psi})$ and $p(0)=\tfrac14\big(2
+i\braket{\psi|U|\psi}-i\overline{\braket{\psi|U|\psi}}\big)
=\tfrac12(1+\mathrm{Im}\braket{\psi|U|\psi})$. The real test returns only
$\cos2\pi\varphi$, an even function of $\varphi$, so $\varphi$ and $1-\varphi$
give the same $p(0)$; the imaginary test returns $\sin2\pi\varphi$, whose sign
breaks the ambiguity.

\textbf{Exercise 3.} After the controlled powers the ancilla holds
$\tfrac{1}{\sqrt8}\sum_{j=0}^{7}e^{i2\pi(3/8)j}\ket{j}=\QFT\ket{3}$ since
$\varphi=3/8=k/2^{t}$ with $k=3$. The inverse QFT returns $\ket{3}=\ket{011}$
and the measurement reads $011$, i.e.\ $\tilde\varphi=3/8$ with certainty.

\textbf{Exercise 4.} With $T=8$, $\varphi=\tfrac13$. Using
$\abs{\gamma_{k}}^{2}=\tfrac{1}{T^{2}}\frac{\sin^{2}(\pi T(\varphi-k/T))}
{\sin^{2}(\pi(\varphi-k/T))}$: for $k=3$,
$\varphi-k/T=\tfrac13-\tfrac38=-\tfrac{1}{24}$, so
$\abs{\gamma_3}^{2}=\tfrac{1}{64}\frac{\sin^2(\pi/3)}{\sin^2(\pi/24)}
\approx\tfrac{1}{64}\cdot\frac{0.75}{0.01704}\approx0.687$.
For $k=2$, $\varphi-k/T=\tfrac13-\tfrac14=\tfrac{1}{12}$, so
$\abs{\gamma_2}^{2}=\tfrac{1}{64}\frac{\sin^2(2\pi/3)}{\sin^2(\pi/12)}
\approx\tfrac{1}{64}\cdot\frac{0.75}{0.06699}\approx0.175$.
The nearest level $k=3$ ($=0.375$) is most likely, as expected since it is the
closest $t$-bit approximation to $0.333\ldots$

\textbf{Exercise 5.} Take $\epsilon=0.1$, $n=6$. Then
\[
  t=n+\big\lceil\log_2\big(2+\tfrac{1}{2\cdot0.1}\big)\big\rceil
   =6+\lceil\log_2(2+5)\rceil=6+\lceil\log_2 7\rceil=6+3=9
\]
ancilla qubits.

\textbf{Exercise 6.} Kitaev's method spends \emph{circuit depth}: applying
$U^{2^{j}}$ (which costs $2^{j}$ applications of $U$) amplifies the phase so that
each bit is read with high confidence, and the total query cost sums to
$O\big(\sum_j 2^j\big)=O(1/\epsilon)$. The single-qubit estimator uses only $U$
and learns $\varphi$ through a probability that must be resolved to
$O(\epsilon)$; by Bernoulli variance this needs $O(1/\epsilon^2)$
\emph{repetitions}. Depth buys a quadratic saving over samples.

\textbf{Exercise 7.} QPE acts on each eigencomponent independently. With
probability $\abs{c_0}^2=0.8$ the measurement returns $\varphi_0$ exactly (it is
$t$-bit representable) and the system collapses to $\ket{\psi_0}$; with
probability $\abs{c_1}^2=0.2$ it returns $\varphi_1$ and the system collapses to
$\ket{\psi_1}$. The measured phase and the surviving eigenstate are perfectly
correlated: outcome $\varphi_k$ always accompanies $\ket{\psi_k}$.
\end{solution}
